\section{R17 执行记录：DSM/cluster V-split 可行性 probe}

\subsection{本轮目标}

R3--R11 的 V-split 方向把 \(D=128\) 的 value/state 维度拆给多个 CTA，
但代价是每个 CTA 都要读取一份相同的只读 tile。R17 只验证一个更窄的问题：
在 B300 上，cluster DSM 能否让两个 CTA 共享同一个 tile，并在一个简化的
V-split 读模式中比 duplicate load 更快。

新增的 \code{challenge/dsm\_cluster\_probe.cu} 不接入 \FlashKDA{} K2，
也不改变默认生产宏。它只启动 cluster size 为 \((2,1,1)\) 的两个 CTA：
duplicate path 中两个 CTA 都把完整 tile 读入自己的 shared memory；DSM
path 中 rank 0 读入一次，rank 1 通过
\code{cluster.map\_shared\_rank} 读取 rank 0 的 shared memory。两个 CTA
分别消费 tile 的前后半段，并把两半求和结果用于 exactness 检查。

\subsection{实现修正和编译}

编译前对 probe 做了三个最小修正，均局限于
\code{challenge/dsm\_cluster\_probe.cu}：

\begin{itemize}
  \item host 侧 exactness 不再解引用 device pointer，而是按输入初始化公式
        在 host 上重建期望值；
  \item kernel 末尾增加 CTA 内 shared-memory reduction，避免只写
        \code{threadIdx.x==0} 的局部累加；
  \item 最大 tile 为 16,384 个 \code{float}，动态 shared memory 为 64 KiB，
        因此 launch 前显式设置
        \code{cudaFuncAttributeMaxDynamicSharedMemorySize}。
\end{itemize}

本轮首次同步时，远端 checkout 缺少修正后的 probe 源文件，故先确认目标目录
和 SSH 身份，再在获得写权限后重试。最终使用显式密钥和
\code{IdentitiesOnly} 完成同步；两个文件均返回成功，随后才提交 Slurm 作业。
仓库中只记录脱敏后的命令形式：

\begin{lstlisting}
scp -F none -i <SSH_KEY> -o IdentitiesOnly=yes \\
  challenge/dsm_cluster_probe.cu challenge/cluster_launch_smoke.cu \\
  <USER>@<B300_HOST>:<B300_REPO>/assignment02/team/c1_flashkda/challenge/
\end{lstlisting}

这一步解决的是 R17 的源文件同步权限/身份问题，不是放宽集群权限；实际
编译和运行仍通过 B300 的 Slurm GPU allocation 完成。

B300 端环境为 \code{infrab300}、CUDA 13.0，CUTLASS 路径沿用
\code{FlashKDA/cutlass}。编译命令如下：

\begin{lstlisting}
ssh infrab300 'set -euo pipefail
  C1=<B300_REPO>/assignment02/team/c1_flashkda
  cd "$C1/challenge"
  CUTLASS="$C1/FlashKDA/cutlass"
  nvcc -std=c++17 -O3 -arch=sm_103a --expt-relaxed-constexpr \
    -I"$CUTLASS/include" -I"$CUTLASS/tools/util/include" \
    -I"$CUTLASS/examples/common" -I. \
    -o cluster_launch_smoke cluster_launch_smoke.cu
  nvcc -std=c++17 -O3 -arch=sm_103a --expt-relaxed-constexpr \
    -I"$CUTLASS/include" -I"$CUTLASS/tools/util/include" \
    -I"$CUTLASS/examples/common" -I. \
    -o dsm_cluster_probe dsm_cluster_probe.cu'
\end{lstlisting}

两个 binary 分别约为 969 KiB 和 995 KiB。先用 job 25001 做
\code{cudaLaunchKernelExC} smoke test：

\begin{lstlisting}
srun --partition=gpu --gres=gpu:1 --cpus-per-task=4 --time=00:03:00 \
  ./cluster_launch_smoke 2>&1 | tee ../results/r17_cluster_launch_smoke_b300.log
\end{lstlisting}

输出显示普通 kernel 和读取 \code{cooperative\_groups::this\_cluster()} 的
cluster kernel 均为 \code{launch=no error}、\code{sync=no error}。对应
结构化结果为 \code{results/r17\_cluster\_launch\_smoke\_b300.json}。

\subsection{多规模 probe 命令}

正式测量为 job 25002。每个 cluster 数量下 probe 内部固定测试
\(1024,2048,4096,8192,16384\) 个 \code{float}，\code{iters=200}，
\code{repeats=9}；输出的时间已经除以 \code{iters}，为单次 producer/consumer
循环的 CUDA-event median。

\begin{lstlisting}
srun --partition=gpu --gres=gpu:1 --cpus-per-task=4 --time=00:15:00 \
  bash -lc 'set -euo pipefail
    for c in 8 16 32 64; do
      echo "=== clusters=${c} iters=200 repeats=9 ==="
      ./dsm_cluster_probe "$c" 200 9
    done' \
  2>&1 | tee ../results/r17_dsm_cluster_probe_b300.log
\end{lstlisting}

完整 raw log 保存为
\code{results/r17\_dsm\_cluster\_probe\_b300.log}，解析后的 JSON 保存为
\code{results/r17\_dsm\_cluster\_probe\_b300.json}。

\subsection{结果}

所有 20 个规模点的 duplicate 与 DSM exactness 都为 true。下表中的
\(\mathrm{dup}/\mathrm{dsm}\) 小于 1 表示 DSM 比 duplicate 慢。

\begin{center}
\small
\begin{tabular}{@{}r r r r r c c@{}}
\toprule
clusters & elems & duplicate (ms) & DSM (ms) & dup/DSM & dup exact & DSM exact \\
\midrule
8  & 1024  & 0.0023064 & 0.0026262 & 0.878214 & true & true \\
8  & 2048  & 0.0037944 & 0.0044792 & 0.847116 & true & true \\
8  & 4096  & 0.0066885 & 0.0080749 & 0.828307 & true & true \\
8  & 8192  & 0.0126046 & 0.0154166 & 0.817600 & true & true \\
8  & 16384 & 0.0243811 & 0.0300067 & 0.812522 & true & true \\
16 & 1024  & 0.0023093 & 0.0026270 & 0.879043 & true & true \\
16 & 2048  & 0.0037974 & 0.0044776 & 0.848097 & true & true \\
16 & 4096  & 0.0066853 & 0.0080749 & 0.827911 & true & true \\
16 & 8192  & 0.0126032 & 0.0154142 & 0.817633 & true & true \\
16 & 16384 & 0.0243814 & 0.0300142 & 0.812329 & true & true \\
32 & 1024  & 0.0023086 & 0.0026286 & 0.878264 & true & true \\
32 & 2048  & 0.0037976 & 0.0044776 & 0.848133 & true & true \\
32 & 4096  & 0.0066851 & 0.0080730 & 0.828088 & true & true \\
32 & 8192  & 0.0126027 & 0.0154066 & 0.818010 & true & true \\
32 & 16384 & 0.0243795 & 0.0300157 & 0.812226 & true & true \\
64 & 1024  & 0.0023190 & 0.0026389 & 0.878797 & true & true \\
64 & 2048  & 0.0037965 & 0.0044771 & 0.847974 & true & true \\
64 & 4096  & 0.0067645 & 0.0081573 & 0.829257 & true & true \\
64 & 8192  & 0.0127656 & 0.0155750 & 0.819619 & true & true \\
64 & 16384 & 0.0246494 & 0.0302826 & 0.813981 & true & true \\
\bottomrule
\end{tabular}
\end{center}

从 JSON 汇总看，\(\mathrm{dup}/\mathrm{dsm}\) 的范围为
\([0.812226,0.879043]\)。换成 DSM 相对 duplicate 的开销，即
\(\mathrm{dsm}/\mathrm{dup}\)，范围为 \([1.1376,1.2312]\)。在这个模型中，
DSM path 没有把 duplicate load 变成收益，反而慢约 13.8\%--23.1\%。
cluster 数量从 8 到 64 时结果几乎不变，说明本 probe 主要测到的是单个
2-CTA cluster 内的 DSM mapping/synchronization 代价，而不是 grid 规模。

\subsection{局限和接入判断}

这个 probe 的结论只覆盖 DSM 共享 tile 这个 primitive，不能直接外推到完整
\FlashKDA{} K2：

\begin{itemize}
  \item probe 的消费阶段只是线性求和，没有 TMA pipeline、MMA、state
        recurrence、Neumann inverse 或 output epilogue；
  \item duplicate path 在两个 CTA 中各自读 global memory，但没有模拟真实
        K2 中 TMA descriptor、shared-memory ring buffer 和 barrier 的固定成本；
  \item DSM path 需要 \code{cluster.sync()} 和 remote shared-memory mapping，
        本轮没有尝试 overlap producer load 与 consumer compute；
  \item tile 以 contiguous float 表示，没有覆盖 BF16 vectorized load、
        swizzle、bank conflict 和真实 \(D=128\) state layout；
  \item 本轮没有改 \FlashKDA{} extension、没有打开任何默认生产宏，也没有把
        DSM path 接入 benchmark。
\end{itemize}

\begin{decision}{R17 接入决策}
在当前证据下，不把 DSM/cluster V-split 接入生产 K2。B300 支持 cluster
launch，且 DSM path 的 exactness 在 probe 中成立；但 20 个规模点全部显示
DSM median 慢于 duplicate median。若后续继续这个方向，应先设计能隐藏
\code{cluster.sync()} 和 remote DSM load 延迟的真实 TMA/MMA pipeline，
再与 R11/R12 的 V-split 候选在同一端到端 benchmark 下比较。
\end{decision}

\begin{record}{R17 结论}
R17 的 DSM/cluster primitive 在 B300 上可编译、可启动、可验证 exact，但
不是当前 V-split 的直接性能答案。job 25002 中 duplicate 和 DSM 的
exactness 全部为 true；DSM/duplicate median 比值为 1.1376--1.2312，说明
简化模型下 DSM 比 duplicate load 慢约 13.8\%--23.1\%。默认生产 kernel
保持不变，本轮只作为 DSM 方向的负向可行性证据。
\end{record}
